\subsection{The full spherical order-five test has no further obstruction}

\begin{statement}{D23.F39}{Proved}{Sixty independent simple factors and a trivial finite order-five cohomology set}
\end{statement}
Use the literal degree-$6720$ spherical Nielsen action of D23.P29.
Let $\Gamma$ be the image of $M(0,5)$ and $P$ the image of its
pure subgroup $K(0,5)$. Its sixty blocks have size $112$.
With the same local group as in D23.C23, put
\[
 L=\operatorname{PSO}_6^-(3),\qquad
 S=L'=\operatorname{P\Omega}_6^-(3),\qquad |S|=3265920.
\]
Then $\Gamma$ contains the full independent product
\[
                    N=S^{60}\subseteq P'.
\]
The quotient $P/N$ is solvable and has order prime to five.
For $R=(\sigma_4\sigma_3\sigma_2\sigma_1)^2$ in this finite
image and $\alpha=\operatorname{Ad}(R)$, one has
\[
                         H^1(C_5,P)=1.
\]
In particular every finite order-five cocycle is a pure
coboundary. All eighteen prescribed cocycles of P29 have
explicit primitives in $N$; these primitives can each be lifted
to an ordinary commutator word in $K(0,5)$. Thus even requiring
zero integral abelianization of a finite primitive does not
exclude any of the eighteen candidates in this representation.

\begin{proof}
We first prove the full simple-product assertion, rather than
infer independence from separate local images. Label the blocks
$(k,j)$ for $0\leq k<5$, $0\leq j<12$, by transporting the original
twelve blocks with $R^k$. They are numbered $12k+j$.
A breadth-first traversal of the four positive half-twists
chooses literal permutations $T_b$ carrying block zero to block
$b$, with $T_0=1$. Normalize local coordinates by these transports.
For every half-twist $g$ and block $b$, the transition
\[
                    T_{q_g(b)}^{-1}gT_b
\]
restricts to an element of $L$ on block zero. All $240$ such
memberships are checked exactly. Their multiplication rule
therefore embeds $\Gamma$ into $L\wr Q$, where $Q$ is its
sixty-block image. The block-zero stabilizer image is all of
$L$: it is contained in $L$ by this calculation and contains
the old block-stabilizer image from the first three half-twists,
which is $L$ by F30 and C23. The action on the sixty blocks is
transitive, so the same assertion holds on every block.

Exact block calculations give
\[
 |Q|=18807399383040,\qquad
 |Q_P|=156728328192=2^{15}3^{14},
\]
where $Q_P$ is generated by the ten correctly marked pure
braids $A_{ij}$. Their color action is trivial; the four
half-twists give the literal $S_5$ action on the five positions
of the exceptional conjugacy class. Consequently the image
of the pure group is exactly the kernel of this color action.
The quotient $|Q|/|Q_P|=120$ also checks its full permutation
image. The stabilizer of block zero in $Q$ has orbits
\[
 \{0\},\quad \{9,11\},\quad
 \{1,2,3,4,5,6,7,8,10\},\quad \{12,\ldots,59\}.
\]
The three orbits of distinct ordered pairs have representatives
$(0,9)$, $(0,1)$, and $(0,12)$.

Here is one explicit word that separates all three types.
For $g_i=\sigma_i$, use the signed-letter word
\[
 w=(1,-3,-2,-4,-3,1,3,-4,1,-3,2,-4,2,-1,1,3),
\]
where a negative letter denotes the inverse and the displayed
order is ordinary left-to-right group multiplication. Its
block permutation has order twenty. Hence $w^{20}$ belongs to
the block kernel, whose normalized coordinates lie in $L$.
As $L/S=C_2$, the element $u=w^{40}$ belongs to $S^{60}$.
The certificate checks this membership on every block and
finds the local orders
\[
       \operatorname{ord}(u_0)=7,\qquad
       u_1=u_9=u_{12}=1.
\]

Let $J$ be the normal closure of $u$ in $\Gamma$. The product
$S^{60}$ is normal in the ambient wreath product, so
$J\subseteq S^{60}$. It is also contained in $P$, since $u$
fixes all blocks and $P$ is normal. Its block-zero projection
is a nontrivial normal subgroup of $L$ contained in the
nonabelian simple group $S$, and thus is $S$. Transitivity
gives surjectivity onto all sixty copies of $S$.
For any of the three representative pairs, the projection of
$J$ is a subdirect subgroup of $S\times S$. By Goursat's lemma
it is either the full product or the graph of an automorphism.
The displayed unequal local orders rule out a graph.
The three pair orbits therefore imply full projection onto
every distinct pair. The simple-product lemma proved in F30
now gives $J=S^{60}$. This proves $N\subseteq P$; perfectness
of $N$ then gives $N\subseteq P'$.

The kernel of $P/N\to Q_P$ embeds into $(L/S)^{60}=C_2^{60}$.
Moreover each of the five sector restrictions of each pure
block generator belongs to the old order-$216$ block group;
these fifty memberships are checked exactly. That old group
is solvable, with derived orders $216,72,18,9,1$ by F30.
Thus $Q_P$ is a subgroup of a product of five solvable groups,
so is solvable. The preceding extension proves that $P/N$ is
solvable and its order has no factor five.

We give the cohomology argument explicitly. If a finite
solvable group $B$ has order prime to five and $\alpha$ is an
automorphism of order dividing five, every $C_5$ cocycle in
$B$ is a coboundary. For an abelian $B$, use additive notation.
If $\sum_{i=0}^4\alpha^i(c)=0$, then with
$T=\sum_{i=0}^4 i\alpha^i(c)$ one has
$(1-\alpha)T=-5c$. Multiplication by five is invertible on
$B$, so $c=(1-\alpha)(-T/5)$. For a solvable $B$, induct on
its derived length, using a nontrivial abelian last derived
subgroup. First kill the cocycle in the quotient, lift the
chosen conjugator, and then kill the resulting cocycle in
that abelian subgroup by the preceding argument. This proves
the assertion without an unproved coprime-conjugacy input.

Now take any cocycle $f\in P$, with
$\alpha^4(f)\alpha^3(f)\alpha^2(f)\alpha(f)f=1$.
Apply the preceding argument in $P/N$, and choose $p\in P$
lifting a primitive there. The cohomologous cocycle
\[
                          f'=\alpha(p)fp^{-1}
\]
belongs to $N$. The automorphism $\alpha$ permutes the sixty
simple factors in twelve free cycles of length five, because
$R$ carries sector $k$ to sector $k+1$. Also $R^5=1$ literally.
Choose local coordinates by transporting along these cycles;
this removes any intermediate inner coordinate changes.
For one cycle, writing the cocycle as $(c_0,\ldots,c_4)$, its
norm condition is $c_1c_2c_3c_4c_0=1$.
Choose $k_0=1$ and recursively $k_j=k_{j-1}c_j$ for
$j=1,\ldots,4$. The norm condition supplies the closing
equation $k_0=k_4c_0$. Therefore
$c_j=k_{j-1}^{-1}k_j$, exactly the equation
$f'=\alpha(k)^{-1}k$ on that cycle. Do this on all twelve
cycles. It follows that $f=\alpha(kp)^{-1}kp$ in $P$.
This proves $H^1(C_5,P)=1$.

Finally P29 supplies eighteen literal primitives $K$ with
$K(f^{-1}R)K^{-1}=R$. Re-expressing all $60$ local coordinates
in the independently chosen transports above verifies that
each belongs to $S$. Hence these entire permutations lie in
$N\subseteq P'$. The image of the ordinary commutator subgroup
of a group under a finite quotient is the commutator subgroup
of the image: lift each factor of a finite product of
commutators. Thus each primitive has an ordinary commutator
lift in $K(0,5)$, and that lift has zero abelianization.
This is a finite compatibility statement, not a lift of an
eighteen-pair candidate to an actual full-GT element.
\end{proof}

\paragraph{Exact source normalization and abelian constraint.}
Lochak--Schneps, \emph{Inventiones Mathematicae} 127 (1997),
Theorem 2, printed p.~573, gives, at multiplier one,
$f(X_1,X_2)=\alpha(k)^{-1}k$ with $k\in\widehat{K(0,5)}$.
Their p.~572 fixes $R=(\sigma_4\sigma_3\sigma_2\sigma_1)^2$
and $\alpha(X_i)=X_{i+3}$, exactly P29's checked marking.
Equivalently $k(f^{-1}R)k^{-1}=R$; the inverse and side of
the conjugator are thereby fixed. The five-generator
presentation on printed p.~578 has all relators of exponent
sum zero in each generator, so its abelianization is
$\mathbb Z^5$, and the profinite abelianization is
$\widehat{\mathbb Z}^5$. Since $f$ is a commutator word,
$k^{\rm ab}$ is fixed by the cyclic permutation of the five
basis vectors. It therefore has the exact form
$t(1,1,1,1,1)$, with $t\in\widehat{\mathbb Z}$.
Unlike H10, this does not force $t=0$; the eighteen finite
primitives above admit even that stronger normalization.

\paragraph{Certificate and external inputs.}
\path{agents/five_level_six/hurwitz_spherical_socle.py}
reconstructs all sixty transports, checks all $240$ transitions,
both block orders, the three pair orbits, all fifty sector
memberships, the explicit word's order witnesses, and the
eighteen primitives on all sixty blocks. The JSON and NPZ
have the same stem. D23.C23 supplies the exact identification
and nonabelian simplicity of $S$; D23.F30 supplies the old local
group and its solvable block quotient. D23.P29 supplies the
literal spherical representation and the eighteen cocycles.
No unverified full degree-$6720$ group order or maximal
solvable quotient identification is used.
